% =====================================================================
%  RESUMO
% =====================================================================

\begin{resumo}

O ciclismo virtual combina rolos de treino inteligentes, simulação de percursos
e análise fisiológica, mas o fluxo de trabalho do ciclista costuma ser
fragmentado entre aplicativos de rota, treino e análise. Este trabalho
apresenta o \argus{}, um ecossistema de software multiplataforma e de código
aberto para simulação, exploração e análise de percursos reais no ciclismo
virtual. A tese que orienta o projeto é que o mesmo fluxo atende a quem vai
disputar uma prova e quer reconhecer o percurso, a quem quer conhecer e pedalar
um lugar que talvez nunca visite presencialmente e a quem quer treinar com
dados fisiológicos. O GPX entra, o rolo simula a estrada, a telemetria é
analisada e, no \textit{desktop}, o histórico permanece no dispositivo; no
Android, a sessão ao vivo cobre a importação, a simulação e a telemetria, sem
FIT ou histórico. O \textit{desktop} foi desenvolvido em Go, com interface
JavaScript embutida por meio do Wails; a mesma interface é reutilizada no
Android por meio do Capacitor, onde a física e o GPX são reimplementados em
JavaScript. O projeto adota princípios de \textit{Clean Architecture}, os
padrões \textit{Facade} e \textit{Repository} e persistência local em SQLite no
\textit{desktop}. A simulação física do \textit{desktop} utiliza a equação de
Martin para estimar a velocidade a partir da inclinação; o módulo de análise do
\textit{desktop} calcula Potência Normalizada (NP), Fator de Intensidade (IF),
\textit{Training Stress Score} (TSS), \textit{Training Impulse} (TRIMP),
\textit{Performance Management Chart} (PMC), recuperação da frequência cardíaca
(HRR) e desacoplamento cardiovascular, com fórmulas e limitações documentadas.
O trabalho adota a \textit{Design Science Research} e avalia o artefato por meio
de um estudo piloto com o próprio autor, complementado por testes de integração
e de exportação de arquivos \texttt{.FIT}. Foram validados o \textit{desktop}
Windows/Linux no fluxo completo e o telefone Android em sessão ao vivo;
iOS/iPadOS, tablet Android e \textit{desktop} macOS permanecem como preparação
arquitetural ou layout projetado, sem validação em dispositivo. O código é distribuído sob GPLv3 e o registro da
versão inicial foi depositado no Instituto Nacional da Propriedade Industrial
(INPI, Processo BR512026003786-7).

\textbf{Palavras-chave}: Ciclismo virtual. Percursos reais. Código aberto.
Análise fisiológica. Multiplataforma. \textit{Design Science Research}.

\end{resumo}
